\begin{lemma}
The incident core preserves edge adjacency and proper coloring, both on the
full original edge type and after retaining an arbitrary set $F\subseteq E$.
Precisely, two original edge identities are adjacent in $G^{\circ}$ if and
only if they are adjacent in $G$; consequently, every partial coloring of
$E$ is proper on $G^{\circ}$ if and only if it is proper on $G$.  Moreover,
for every $F\subseteq E$, the analogous two equivalences hold between
$G^{\circ}|F$ and $G|F$, using the identical retained edge subtype.
\end{lemma}
\begin{proof}
Fix edge identities $e$ and $f$.  If they are adjacent in $G^{\circ}$,
they have a common core endpoint $x$.  Applying
\noderef{IncidentCoreIncidence} twice shows that the underlying vertex of
$x$ is incident with both $e$ and $f$ in $G$, so the identities are adjacent
there.  Conversely, if a vertex $v$ is incident with both identities in
$G$, regard $v$ as an incident vertex using $e$ as its witness.  The reverse
direction of \noderef{IncidentCoreIncidence}, first for $e$ and then for
$f$, makes this core vertex a common endpoint.  This proves full adjacency
equivalence.

The definition \noderef{IsProperEdgeColoring} tests only distinct edge
identities, adjacency, and their assigned colors.  The edge type and coloring
function are unchanged, and the preceding paragraph preserves adjacency;
substitution in that definition proves properness in both directions.

Now fix $F\subseteq E$.  By \noderef{RestrictEdges}, an edge of either
restriction is the same pair consisting of an identity $e\in E$ and a proof
that $e\in F$, and its endpoint pair is inherited without change.  Repeating
the common-endpoint argument above, while forgetting only the membership
proof in $F$, proves adjacency equivalence for $G^{\circ}|F$ and $G|F$.
Substitution into \noderef{IsProperEdgeColoring} again proves the restricted
properness equivalence.  At no point are two retained, parallel identities
identified.
\end{proof}
